%%%%% 3. Project management %%%%%
\chapter{Project management}\label{ch:project-management}

\section{Agile}\label{sec:agile}

Agile is a popular methodology within software development and wider business
contexts \citep{anifa2024systematic}. Whereas other methodologies set out a
pre-defined timetable and structure with deadlines and aims attached to each
step, the agile methodology is more flexible and accounts for the fact that most
software, and in this case research, projects cannot be defined with complete
accuracy at the start of the project. The founding principles of agile treat
this as a design goal rather than a concession, and ask that changing
requirements be welcomed even late in development \citep{beck2001manifesto}. The
systematic review by \citet{dyba2008empirical} finds empirical support for the
claim: by planning in detail only the features of the current cycle, teams were
able to absorb later changes in requirements at a lower cost.

The agile methodology was followed in the following steps:
\begin{itemize}
    \item \textbf{Defining the goals.} The aim of this agile method was defined
    as a SMART goal \citep{doran1981smart}. The aim was to create a complete
    thesis document by the 25th of August 2026. The thesis needed to be a
    self-standing and complete document describing a new methodology of creating
    differentially private shared memories. This aim was specific, measurable,
    assignable, realistic and time-related. Having this aim clearly in mind for
    each weekly sprint plan allowed the author to keep a clear focus and
    prioritise where needed, for example when deciding between code expansion
    and writing up the thesis, where the latter choice was made.

    \item \textbf{Product backlog.} The product backlog was initially created
    from the work packages set out in the proposal. While these were, in
    particular at the beginning, helpful in defining a structure, later weeks
    diverged from the exact set of tasks to better suit the requirements of the
    project. The backlog is defined as an emergent list rather than a fixed one
    \citep{schwaber2020scrum}, so this divergence is the intended behaviour of
    the framework rather than a departure from it.

    \item \textbf{Sprints.} Each week marked the beginning of a new sprint,
    where the author reviewed the items on the Kanban board and decided which
    items to carry forward, which new items to add, and which were finished.

    \item \textbf{User stories.} The tasks for this project were all grouped
    under a user story. The stories for this project are listed here:
    \begin{itemize}
        \item \textit{``As a Researcher, I want to assemble all proposal
        sections into a polished, well-formatted document.''}
        \item \textit{``As a Researcher, I want to define concrete evaluation
        metrics and experimental protocols.''}
        \item \textit{``As a Researcher, I want to refine my technical
        architecture and verify tool feasibility.''}
        \item \textit{``As a Researcher, I want to create a realistic project
        plan with work breakdown structure, timeline, and risk mitigation
        strategies.''}
        \item \textit{``As a Researcher, I want to define a precise, scoped
        research question grounded in essential literature.''}
        \item \textit{``As a Researcher, I want to understand differential
        privacy.''}
        \item \textit{``As a Developer, I want to have a clean repo with clear
        structure.''}
        \item \textit{``As a Developer, I want to create a complete architecture
        in which to run my agent.''}
        \item \textit{``As a Researcher, I want to specify a detailed
        methodology with a formally defined privacy mechanism.''}
    \end{itemize}
    These stories were all completed with this thesis, and provided a guiding
    direction for the weekly sprint planning.

    \item \textbf{Refinement.} This was done at the end of every week before
    planning the sprint for the next week, where the author decided whether the
    timelines for any item had to be extended or changed, and whether they were
    still relevant.

    \item \textbf{Feedback.} Feedback was provided by the author's supervisor,
    Luca Citi, with whom the author met at a bi-weekly, and later in the process
    weekly, pace. In addition to receiving structured feedback in each session,
    these meetings allowed the author to align on priorities, troubleshoot, and
    discuss risks and impediments. The feedback is recorded alongside the sprint
    items on the Kanban board.
\end{itemize}

Two elements of the framework were not adhered to. The first is the daily
scrum. Since the author worked by herself, a daily status meeting had no
audience and was dropped. The second is the retrospective, which was not
followed as a written ceremony but as an internal process on a continuous basis,
where the author assessed what worked in her working process and what did not.
Since she is the sole contributor to the paper, she did not need to write these
out. This thesis therefore implements parts of the agile methodology rather than
the framework in full. Doing so is the common case rather than the exception, as
\citet{fitzgerald2006customising} show in a three-year study at Intel Shannon,
where the agile methods in use were tailored to the practices of the team rather
than adopted wholesale.

Risk was handled at the weekly refinement rather than through a separate risk
register. Each week the author reviewed which items were at risk of slipping and
decided whether to extend, re-scope or drop them, and the supervisor meetings
surfaced the risks that were not visible from the board alone. The pivots that
followed, including the move from a cloud instance to local hardware and the
change of hypotheses, are set out in Chapter~\ref{ch:interim-results}.

The agile methodology has allowed the author to re-adjust her targets based on
what worked and what did not work. The entire
Chapter~\ref{ch:interim-results} speaks of the numerous pivots that were
required to get this thesis to its current state. The author is therefore
content with the choice of methodology.

\section{Tool use}\label{sec:tool-use}

\subsection{Tuleap}\label{sec:tuleap}

Tuleap is the university-provided agile methodology tracking tool. All weekly
feedback, user stories, and tasks were loaded here. The Kanban board it provides
allowed the author to organise the weekly sprint planning cleanly, and to move
items across the board. A weekly screenshot of a board is included below.

\begin{figure}
    \centering
    \includegraphics[width=1\linewidth]{Tuleap.png}
    \caption{Weekly Kanban board excerpt from Tuleap}
    \label{fig:Tuleap}
\end{figure}

\subsection{Overleaf}\label{sec:overleaf}

Overleaf was used as the interface for writing this thesis in \LaTeX. The
website simplifies compiling the code, and by linking it to GitHub, the source
can still be copied to a local repository for offline work. The interface also
allows concurrent views of the table of contents, the compiled PDF, and the
source, which the author preferred over working in the source alone.

\subsection{GitLab and GitHub}\label{sec:gitlab}

This author used both GitLab and GitHub as repositories. The GitLab repository,
hosted on the university's \texttt{cseegit} instance, is the marked repository
and is where all changes, merges, and edits took place. The GitHub repository is
a public mirror, kept so that the final repository could be added to the
author's public profile. The two are not maintained by hand: the \texttt{origin}
remote is configured with one fetch URL and two push URLs, so a single
\texttt{git push origin} updates both at once.

The author implemented a clear structure of creating a new branch for each task,
and followed a consistent naming convention of
\texttt{<task-number>\_<name>}, for example
\texttt{21589\_tweak\_code}. After adding and committing any local code, it was
pushed to GitLab. There, the author would conduct a final review and merge the
branch into \texttt{master} using the pre-structured merge request template shown
in Figure~\ref{fig:Gitlab_merge_template}. The template is again part of the
author's general tech stack and was created as part of the group project module
at the University of Essex. It sets out a summary field, a checklist covering
tests, code quality, dependencies and documentation, and a block into which the
output of the passing test suite is pasted, so that every merge carries its own
evidence.

\begin{figure}
    \centering
    \includegraphics[width=1\linewidth]{Gitlab_merge_template.png}
    \caption{GitLab merge request template}
    \label{fig:Gitlab_merge_template}
\end{figure}

\subsection{Code quality tooling}\label{sec:code-quality-tooling}

The checks below are run together with a single \texttt{make ci} command, which
runs Black in check mode, Ruff, mypy and pytest in sequence and stops at the
first failure. The same command is named in the merge request checklist, so no
branch reaches \texttt{master} without it having passed.

\begin{itemize}
    \item \textbf{Ruff.} Ruff is the linter. It is configured to check
    pycodestyle errors and warnings, pyflakes rules, import ordering, PEP 8
    naming, outdated syntax, and the bugbear rule set, across \texttt{src/},
    \texttt{evaluation/}, \texttt{scripts/} and \texttt{tests/}. Line-length
    errors are ignored, since Black already enforces the 100-character limit.

    \item \textbf{Black.} Black is the formatter. It rewrites the same four
    directories to a fixed style at a line length of 100, which removes
    formatting from the set of things a review has to consider.

    \item \textbf{mypy.} mypy is the static type checker. It runs over
    \texttt{src/agent\_memories/} and \texttt{evaluation/} and rejects functions
    with missing or incomplete annotations, as well as returns that silently
    degrade to \texttt{Any}. Third-party libraries that ship no type stubs, such
    as LangGraph, BrowserGym and Playwright, are exempted from the import check
    so that the strictness applies to first-party code alone.

    \item \textbf{pytest.} pytest runs the test suite, which holds 97 test
    functions across 13 files, split into 12 unit files and one integration
    test. Tests are marked as \texttt{unit}, \texttt{integration} or
    \texttt{slow}, and unregistered markers are rejected so that a typo in a
    marker fails the run rather than silently skipping a test. Line and branch
    coverage over \texttt{src/agent\_memories} and \texttt{evaluation} is
    reported on every run, but no minimum threshold is enforced.

    \item \textbf{\texttt{pyproject.toml}.} Rather than use a
    \texttt{requirements.txt} file, this thesis uses a \texttt{pyproject.toml}
    file. The first benefit is that the author was able to differentiate between
    the packages needed to run the project itself and the packages needed only
    to develop it, which are declared as an optional \texttt{dev} group holding
    pytest, Black, Ruff and mypy. The second is that the same file holds the
    configuration for all four of those tools, so a rule change is made in one
    place rather than in four separate configuration files. The third is that
    the project is installed in editable mode from it, which places
    \texttt{src/agent\_memories} on the import path and lets the tests import
    the package directly rather than manipulating \texttt{sys.path}.

    \item \textbf{Helper functions.} The codebase included a Makefile for eight
    simple commands. These were based on previous work the author had conducted
    as part of the group project module at the University of Essex, and have
    become part of her every work cycle since. These included:
    \begin{itemize}
        \item \texttt{make help} Prints a reminder of what each make command
        does, generated from the Makefile's own comments so that it cannot fall
        out of date.
        \item \texttt{make venv} Creates the project virtual environment at
        \texttt{venv/} in the repository root, if it does not exist yet.
        \item \texttt{make install} Installs the project and its development
        dependencies, together with the Chromium binary that Playwright drives
        and the \texttt{punkt\_tab} tokeniser data, both of which are needed for
        the WebArena and BrowserGym run.
        \item \texttt{make format} Auto-runs Black to format the source code in
        the \texttt{src/}, \texttt{evaluation/}, \texttt{scripts/} and
        \texttt{tests/} directories.
        \item \texttt{make lint} Finds and fixes the naming, style and imports
        of the source code with Ruff across the same four directories.
        \item \texttt{make typecheck} Runs mypy over
        \texttt{src/agent\_memories/} and \texttt{evaluation/}.
        \item \texttt{make test} Runs the full test suite with a coverage
        report that names the uncovered lines.
        \item \texttt{make ci} Runs all four checks in their non-fixing form,
        so that the command reports what is wrong rather than quietly repairing
        it.
    \end{itemize}

    \item \textbf{Coding standards document.} The coding standards are the
    common coding standards this author implements in all her projects, and
    includes as a generic document to reference. It originated from the same
    group project module the author had previously partaken in at the University
    of Essex. It was adapted to this project, and has been added to the annex.
    The generic sections cover branch naming, commit messages, commit squashing,
    package structure and dependency management, while the project-specific
    sections cover the remote configuration and the local model setup that
    \texttt{make install} does not handle.

    \item \textbf{Pre-push hooks.} A single pre-push hook fires the two LLM
    calls for the \texttt{decision-curator} and \texttt{thesis-tasks-curator}
    commands described in Section~\ref{sec:llms}. The hook runs them in the
    background and always exits successfully, so a slow or failing curator never
    blocks a push, and it exits silently if the \texttt{claude} binary is not
    installed. The restriction to the primary branch is not enforced by the hook
    itself but inside the \texttt{thesis-tasks-curator} command, which exits on
    any branch other than \texttt{master}. Because the hook lives in a tracked
    \texttt{.githooks/} directory rather than in \texttt{.git/hooks/}, it is
    version-controlled alongside the code and is activated on a fresh clone by
    pointing Git at that directory.
\end{itemize}

\subsection{LLMs}\label{sec:llms}

LLMs were used in this thesis as tools to support the author's progress. While
LLMs can be used to do more comprehensive work, this often results in
high-level, uncreative, and possibly false results. The author therefore used
her subscription to Anthropic's Claude as a tool, without giving up the ultimate
control of the project. The support was divided into the following streams:

\begin{itemize}
    \item \textbf{Logging.} The author created multiple skills to track the work
    that was being done between each commit, and to log it for future purposes.
    The author made a clear difference between skills and commands. With Claude,
    skills can be triggered by anything written in a chat, whereas commands need
    to be invoked specifically, otherwise they will not run. The prompts for the
    different skills and commands used are all attached to the annex.
    \begin{itemize}
        \item First, when pushing to the origin, a \texttt{decision-curator}
        command assesses each decision that was made within the last commit
        messages, and adds in reasoning, where available in the code, to a
        growing memory file. This ensured the author automatically grew a log of
        every decision and change made throughout the code. This decision
        curator was triggered with a pre-push git hook, so the author never had
        to remember it.
        \item Second, when pushing to origin, a \texttt{thesis-tasks-curator}
        command assessed which thesis tasks, as listed in the original work
        breakdown structure in the proposal document, had been completed, and
        how. This allowed the author to have a clean overview of what had been
        done and what still needed to be done. This command only ran on the
        \texttt{master} branch, so that it would only keep track of tasks once
        they were truly complete and merged.
        \item Third, the author added a \texttt{thesis-tasks} skill, which she
        could invoke to check, update, and add to the outstanding tasks list,
        which was originally based on the work breakdown structure.
    \end{itemize}

    \item \textbf{Refining text.} The author used a \texttt{writing-editing-text}
    skill to edit her drafts to sound cleaner and more concise. Importantly, this
    was only used to make the text sound cleaner, and was always based on an
    existing draft of her own, not to write new sections. The main idea is that
    she was still the author, with a smart editor giving feedback and checking
    her writing.

    \item \textbf{Code checks.} On coding, the author used Claude Code to check
    for bugs in her code prior to pushing to origin. As with writing, Claude was
    used to ensure the code stayed consistent, in line with the coding standards
    outlined in \texttt{CODING\_STANDARDS.md} (see annex), and bug-free. Claude
    Code was also used as a discussion partner on why certain things may not have
    been working, and on how to best move forward, on days without a supervisor
    meeting. It also helped provide a check that all of the test cases were
    comprehensive.

    \item \textbf{Source finding.} Lastly, the author used Claude Code to check
    her academic sources, and ensure she did not forget any key sources for a
    topic. She would then use that information to find the referenced source, and
    check whether it should indeed be added to this thesis, and if so, in what
    manner.
\end{itemize}
